\subsection{Comparison across families}
\label{sec:comparison}

\begin{table*}[!tb]
\centering\footnotesize\setlength\tabcolsep{3pt}\renewcommand{\arraystretch}{1.12}
\caption{What each method family shows, what it does not, and how far this corpus has taken it.}
\label{tab:families}
\begin{tabular}{@{}L{2.2cm}L{3.0cm}L{2.5cm}L{3.0cm}L{3.1cm}L{1.4cm}@{}}
\toprule
Family & Shows on its own & Does not show & A stronger reading needs & In this corpus & Cost \\
\midrule
\textbf{Structural and geometric} (task-linked)$^{\dagger}$ & the shape of the space: gap, alignment, effective rank, similarity & what any direction means; no concept & comparison model; random-initialised or shuffled baseline & comparison model \Vcomp{}/\Vcompn{}; calibration \Vncalgeo{}/\Vncalgeon{} & low \\
\addlinespace[3pt]
\textbf{Supervised concept mapping and decodability} & a concept can be read out of $z_\lambda$ by a decoder (Eq.~\ref{eq:probe}) & that the model uses it & permuted labels; random encoder; control task; fixed decoder class & permutation null \Vpermdec{}/\Vpermdecn{}; random encoder \Vlrndec{}/\Vlrndecn{}; selectivity not yet measurable & low \\
\addlinespace[3pt]
\textbf{Feature decomposition and discovery} & a unit's activation tracks the named content (Eq.~\ref{eq:sae}) & that the unit is stable, complete or used & validated names; cross-seed matching; reconstruction & semantic validation \Vsem{}/\Vsemn{}; cross-seed identity none & high \\
\addlinespace[3pt]
\textbf{Attribution, localisation and tracing} & where information is decodable or attended ($\ell$, $p$, $h$) & that the location is used & occlusion or patching confirmation; random-component comparison & input corroboration \Vinp{}/\Vinpn{}; component perturbation \Vinttr{}/\Vinttrn{} & moderate \\
\addlinespace[3pt]
\textbf{Representation intervention and control} & the output is sensitive to $z_\lambda$ (Eq.~\ref{eq:effect}) & that the named concept, not a correlated direction, caused the change & matched control directions; equal-size random units; sham transformations & matched control \Vspca{}--\Vspc{}/\Vspcn{}; no dose--response test & mod.--high \\
\midrule
\textbf{Intrinsic and concept-aligned designs} (tag) & the label depends on $z_\lambda$ only through the concept layer & that nothing else leaks through & leakage test; random concept projections; concept-layer intervention & concept-layer control \Vintrany{}/\Vintrn{}; leakage rarely measured & low--mod. \\
\bottomrule
\end{tabular}
\\[3pt]
\begin{minipage}{\textwidth}\footnotesize
No test in the fourth column is a prerequisite for the family's base reading; a family that lacks one is
narrower, not refuted. Counts are $k/n$ over the family's own records, and the cross-family counts are in
the supplement. \emph{Cost} is what the family needs beyond a forward pass. Which loci each family has
reached is Fig.~\ref{fig:corpus}b, and the sub-types of decoding and intervention are reviewed in
\S\ref{sec:decoding} and \S\ref{sec:intervention}. $^{\dagger}$Task-linked geometry only
(\S\ref{sec:method}), so this row's counts are not comparable with the others; the \Nintrinsic{}
intrinsic-design records are counted in their analytic family and listed again in the tag row.
\end{minipage}
\end{table*}

Table~\ref{tab:families} sets out what each family reports, what its artefact licenses and how far this
corpus has taken it. Decoding is by far the largest family and, with structural analysis and tracing, the
oldest, first appearing in 2022; discovery is two years old and intervention three. Two asymmetries
follow from the table. Intervention gives the claims of greatest practical interest at a high cost and
with the smallest controlled fraction, while the cheapest family addresses availability, which is often
mistaken for use. Readability and validity also pull apart: discovery and intrinsic designs return the
artefacts a clinician can read, yet they rest on an unvalidated naming or leakage step.

\textbf{Combining families.} Two pipelines recur in the medical corpus and a third exists so far only as
a general-domain template. \emph{Discover, then intervene} uses residual-stream
dictionary directions for steering or
patching~\cite{bouzid2025insights,nooralahzadeh2026universal,sadanandan2026mechanistically}, and carries
the caveats of \S\ref{sec:discovery} into the intervention. \emph{Trace, then intervene} exploits an
observational localisation through a training-time change, so far only in a general CLIP text tower on
radiology captions~\cite{abbasi2026layerspecific}. \emph{Decode, then erase, residualise or ablate}
probes for availability and then tests whether a downstream prediction depends on the
concept~\cite{kim2026encoding,nooralahzadeh2026sparse,weber2023posthoc}; it answers the shortcut
question of \S\ref{sec:applications} most directly, and its two core instances that test a downstream
dependence point in opposite directions, residualisation leaving a refitted finding probe
unchanged~\cite{kim2026encoding}$^{\dagger}$ while channel knockout collapses a finding's
score~\cite{nooralahzadeh2026sparse}. Reading one family's result as an answer to another family's
question is the commonest inferential error in this literature, and the evidence-type tag makes it
visible.

